
\documentclass[11pt]{article}
\usepackage{amsmath,amssymb,amsfonts,bm}
\usepackage{geometry,hyperref}
\geometry{margin=1in}

\title{\textbf{SAT.4D Framework --- Phase~IV Summary}\\Cosmology, Perturbations, and Dark Sector}
\author{Nathan McKnight \& Collaborators}
\date{\today}

\begin{document}
\maketitle

\begin{abstract}
Phase~IV extends the SAT.4D framework to cosmology. We derive Friedmann-like
expansion from bundle strain energy, explain the origin of structure via topological
fluctuations, and identify the dark sector with uncoupled or topologically frozen
filament bundles.
\end{abstract}

\tableofcontents

%-------------------------------------------------
\section{SAT Cosmology and Expansion}

Assume a statistically homogeneous and isotropic filament ensemble \( F(x) \).  
The emergent metric is FLRW:
\[
ds^2 = -dt^2 + a(t)^2 d\vec{x}^2.
\]
Each filament contributes strain energy:
\[
\mathcal E[\gamma] = \frac{T}{2} \int d\lambda\, \left\| \frac{d^2\gamma^\mu}{d\lambda^2} \right\|^2,
\]
and the total energy density is
\[
\rho(t) = \rho_b(t) \cdot \varepsilon_b(t).
\]

The expansion rate becomes:
\[
\left(\frac{\dot a}{a}\right)^2 = \frac{\kappa}{3} \rho(t).
\]

Depending on how \( \varepsilon_b \) evolves, SAT strain can reproduce matter, radiation, and dark energy:
\[
\rho(t) = \rho_{\text{M}}\,a^{-3} + \rho_{\text{R}}\,a^{-4} + \rho_\Lambda.
\]

\subsection{Cosmological Constant}
Residual strain in topologically frozen bundles sets:
\[
\rho_\Lambda = \lim_{t\to\infty} \langle \varepsilon_b^{\text{frozen}} \rangle.
\]

\subsection{Inflation}
Initial high strain (\( \varepsilon_b \gg 1 \)) induces rapid expansion.
Relaxation of bundle strain ends inflation naturally.

%-------------------------------------------------
\section{Structure Formation and Perturbations}

\subsection{Sources}
Scalar and tensor perturbations arise from:
\begin{itemize}
  \item Density contrast \( \delta(x,t) = \delta\rho / \bar\rho \)
  \item Orientation field shear \( \delta\tilde g_{ij}^{\text{TT}} \)
  \item Phase vibration modes \( \xi^\mu(\lambda) \)
\end{itemize}

\subsection{Power Spectra}
Filament fluctuations during inflation freeze into curvature perturbations:
\[
P(k) \sim k^{n_s - 1}, \quad n_s \approx 0.96.
\]
Tensor spectrum suppressed unless tension is extreme.

\subsection{Non-Gaussianity}
Topological linking (e.g., triple-bound filaments) induces weak non-Gaussian signals.
Possible topological defects yield CMB anomalies and structure irregularities.

%-------------------------------------------------
\section{Dark Matter and Dark Energy}

\subsection{Sector Classification}

\begin{center}
\begin{tabular}{|c|c|c|c|c|}
\hline
Sector & Q & Phase State & Coupling & Behavior \\\hline
Visible & 1--3 & Locked & Electroweak/color & SM matter \\
Dark Matter & $\geq$1 & Incoherent & Gravity only & Cold, pressureless \\
Dark Energy & $\geq$3 & Locked + frozen & None & $\rho=\text{const}$ \\\hline
\end{tabular}
\end{center}

\subsection{Dark Matter}
Uncoupled filament bundles with residual strain behave as CDM:
\begin{itemize}
  \item Stable under SAT topology
  \item Do not radiate or interact electromagnetically
  \item Gravitate normally and cluster
\end{itemize}

\subsection{Dark Energy}
High-Q frozen bundles contribute constant residual strain:
\[
\rho_\Lambda = \text{Topologically locked strain energy}.
\]

\subsection{Predictions}
\begin{itemize}
  \item No new particles (WIMPs, axions) predicted
  \item No fifth forces or long-range DM interactions
  \item DE equation of state \( w \approx -1 \), possibly evolving slightly
\end{itemize}

%-------------------------------------------------
\section{Conclusion}
Phase~IV demonstrates that SAT.4D can reproduce large-scale cosmic expansion, 
explain the emergence of structure, and account for the dark sector via purely topological
principles. These results arise with no extra fields or fine-tuning, and yield novel observational predictions.

\end{document}
